\documentclass[11pt]{article}
\usepackage[margin=1in]{geometry}
\usepackage{amsmath,amssymb,amsthm,booktabs}
\usepackage[hidelinks]{hyperref}
\setlength{\emergencystretch}{2em}
\newtheorem{proposition}{Proposition}
\newtheorem{corollary}{Corollary}
\newcommand{\norm}[1]{\left\lVert#1\right\rVert}
\newcommand{\R}{\mathbb R}
\title{CASS: Corrected Support Conditions and Operator Geometry}
\author{Technical companion to the author responses}
\date{27 September 2026}
\begin{document}
\maketitle

This companion gives complete proofs for the weighted objective and the actual
hybrid intervention. It distinguishes the submitted operator from the additional
contextual variant. These are conditional mathematical statements; the empirical
comparison of operators and the geometry checks are separate evidence.

\section{The weighted support condition}

Let $U_t\in\R^{d\times r_t}$ have orthonormal columns, let
$A=[U_1\ \cdots\ U_T]$, and write the implemented objective as
\begin{equation}\label{eq:objective}
 \min_c\ \frac12\norm{z-Ac}_2^2+\lambda\sum_{t=1}^T q_t\norm{c_t}_2,
 \qquad q_t=\sqrt{r_t}>0.
\end{equation}
Orthonormality holds within a block; distinct blocks need not be orthogonal.
In the multi-layer construction, the joint blocks are block diagonal and $r_t$
is their total number of columns.

\paragraph{Why the weights matter.}
An independently constructed counterexample makes the missing condition
explicit. In $\R^5$, set
\[
 U_1=[e_1,e_2,e_3,e_4],\qquad U_2=0.6e_1+0.8e_5,
 \qquad z=1.01e_1,\quad\lambda=0.01.
\]
The true support is $S=\{1\}$, with coefficient $(1,0,0,0)$ and noise
$0.01e_1$. Maximum block coherence is $0.6$, so the unweighted condition
$|S|<(1+1/0.6)/2$ holds. However, $q_1=2$ and $q_2=1$. The optimum restricted
to $S$ is $(0.99,0,0,0)$ and has residual $0.02e_1$. Thus
\[
 \norm{U_2^\top(z-U_1\widetilde c_1)}_2=0.012>\lambda q_2=0.01.
\]
The inactive-block KKT condition fails. The full optimum has nonzero second
block (numerically, $\widehat c_{1,1}=0.988125$ and $\widehat c_2=0.003125$).
The same failure occurs for every positive $\lambda<(1+\epsilon)/2$ when
$z=(1+\epsilon)e_1$: changing only a constant relating $\lambda$ to small
noise does not repair the missing weights.

\newpage
\begin{proposition}[Weighted deterministic support bound]\label{prop:weighted}
Suppose $z=A_Sc_S^0+e$, $\norm{e}_2\leq\epsilon$, and $A_S$ has full column
rank. Define
\begin{align*}
 G&=A_S^\top A_S,& P&=A_SG^{-1}A_S^\top,\\
 W_S&=\operatorname{blockdiag}(q_tI_{r_t}:t\in S),\\
 \theta_S&=\max_{j\notin S}\frac1{q_j}
   \sup_{\max_{t\in S}\norm{v_t}_2\leq1}
     \norm{U_j^\top A_SG^{-1}W_Sv}_2,\\
 \eta_S(e)&=\max_{j\notin S}\frac{\norm{U_j^\top(I-P)e}_2}{q_j}.
\end{align*}
If $\theta_S<1$ and $\lambda>\eta_S(e)/(1-\theta_S)$, the unique optimum of
\eqref{eq:objective} has support contained in $S$, and
\begin{equation}\label{eq:error}
 \norm{\widehat c_S-c_S^0}_2
 \leq \frac{\epsilon}{\sigma_{\min}(A_S)}
 +\frac{\lambda\sqrt{\sum_{t\in S}q_t^2}}{\sigma_{\min}(A_S)^2}.
\end{equation}
Exact support follows if every true block norm exceeds the right side of
\eqref{eq:error}. An $O(\epsilon)$ error conclusion additionally requires
$\lambda=O(\epsilon)$ for a fixed dictionary.
\end{proposition}

\begin{proof}
Let $\widetilde c$ minimize the objective restricted to $S$. Full column rank
makes the restricted objective strictly convex, so its minimizer is unique.
The active KKT equations give a block subgradient $v$, with
$\norm{v_t}_2\leq1$, such that
\[
 G(\widetilde c-c_S^0)=A_S^\top e-\lambda W_Sv.
\]
Its residual therefore satisfies
\[
 r=z-A_S\widetilde c=(I-P)e+\lambda A_SG^{-1}W_Sv.
\]
For every $j\notin S$,
\[
 \frac{\norm{U_j^\top r}_2}{q_j}
 \leq\eta_S(e)+\lambda\theta_S<\lambda.
\]
These strict inactive inequalities, together with the restricted KKT equations,
show that extending $\widetilde c$ by zeros gives a global optimum.

All optima have the same fitted value: otherwise strict convexity of squared
loss in the fitted value, together with convexity of the penalty, would make
their midpoint strictly better. Hence all optima have the same residual.
A nonzero inactive block would require equality
$\norm{U_j^\top r}_2=\lambda q_j$, contradicting the strict inequalities.
All optima are thus zero outside $S$, and full column rank gives uniqueness.
Finally,
\[
 \norm{G^{-1}A_S^\top}_2=\sigma_{\min}(A_S)^{-1},\qquad
 \norm{G^{-1}}_2=\sigma_{\min}(A_S)^{-2},\qquad
 \norm{W_Sv}_2\leq\sqrt{\sum_{t\in S}q_t^2}
\]
prove \eqref{eq:error}. Each block error is bounded by the total coefficient
error, giving the stated sufficient minimum-amplitude condition.
\end{proof}

This is a finite KKT argument for an established weighted group-selection
condition, not a claim of new general sparse-recovery theory; see
\href{https://www.jmlr.org/papers/v9/bach08b.html}{Bach (2008), Section 2.8}.
The support-capped finite-iteration solver is not automatically certified to use
a noise-calibrated $\lambda$ satisfying Proposition~\ref{prop:weighted}.

\begin{corollary}[A conservative maximum-coherence condition]
Let $s=|S|$, $\mu=\max_{j\ne t}\norm{U_j^\top U_t}_2$, and
$R_S=\max_{t\in S}q_t/\min_{j\notin S}q_j$. If $(s-1)\mu<1$, then
\[
 \theta_S\leq\frac{s\mu R_S}{1-(s-1)\mu}.
\]
Consequently $\mu<1/(sR_S+s-1)$ implies $\theta_S<1$.
\end{corollary}

\begin{proof}
In the block maximum norm, $\norm{G-I}\leq(s-1)\mu$. The Neumann series bounds
$\norm{G^{-1}}\leq[1-(s-1)\mu]^{-1}$. Moreover,
$\norm{W_Sv}_{\mathrm{block},\infty}\leq\max_{t\in S}q_t$, while the
off-support block row $U_j^\top A_S$ has induced norm at most $s\mu$.
Combining these bounds and dividing by $q_j$ proves the result.
\end{proof}

Replacing $R_S$ by the global maximum-to-minimum weight ratio gives a
support-independent condition. Equal ranks recover $\mu<1/(2s-1)$.
Also $\eta_S(e)\leq\epsilon/\min_{j\notin S}q_j$. A median coherence cannot
replace the maximum. For joint block-diagonal dictionaries, each pair's joint
coherence is the maximum of its layerwise coherences.

\section{Non-vacuity on the submitted dictionaries}

We recomputed the five full dictionaries and all 160 LOTO dictionaries.
The table gives full-dictionary values. ``None'' means the conservative global
weighted corollary certifies no nonempty support, not that recovery is impossible.
\begin{center}
\begin{tabular}{lrrr}
\toprule
Dictionary & Maximum $\mu$ & Global weighted $s$ & Eligible two-block sets\\
\midrule
Llama-3.1-8B &0.994555 &None &10\\
Llama-3.2-3B &0.992069 &1 &1\\
Gemma-2-2B &0.988240 &None &1\\
Qwen3-4B &0.993710 &1 &2\\
Qwen2.5-3B &0.993471 &1 &0\\
\bottomrule
\end{tabular}
\end{center}

The final column checks all $\binom{32}{2}=496$ supports in each dictionary
using a sharper bound for $\theta_S$. For $S=\{s_1,s_2\}$ and $j\notin S$,
write
\[
 [M_1\ M_2]=q_j^{-1}U_j^\top A_SG^{-1}W_S.
\]
Duality and weighted Cauchy--Schwarz give, for any $t\in(0,1)$,
\begin{align*}
 \sup_{\norm{v_1},\norm{v_2}\leq1}\norm{M_1v_1+M_2v_2}_2
 &=\max_{\norm{u}_2=1}\big(\norm{M_1^\top u}_2+\norm{M_2^\top u}_2\big)\\
 &\leq\sqrt{\lambda_{\max}\!\left(
          \frac{M_1M_1^\top}{t}+\frac{M_2M_2^\top}{1-t}\right)}.
\end{align*}
Taking the minimum of evaluated upper bounds remains an upper bound; numerical
scalar optimization need not find the exact supremum for this reasoning.
The maximum bound over inactive blocks, plus a $10^{-6}$ numerical margin,
determines eligibility. Llama-3.1's park-country/word-length pair has bound
$0.9417937318$ before this margin and minimum Gram eigenvalue $0.6363702$.
None of the Llama-3.1 Compound constituent pairs is eligible. Gemma's
country-capital/word-length pair has bound $0.9978110662$.

These geometric conditions do not establish the representation, noise,
regularization or minimum-amplitude assumptions for observed task signatures.
As implementation checks, three planted coefficient/noise instances for each of
the 14 eligible pairs (42 total) recover their support and obey the coefficient
bound under an explicitly admissible $\lambda$. They use the weighted objective
at that prescribed $\lambda$, not the deployment support-cap selection rule.

\section{A direct bound for the submitted hybrid}

Fix the same hidden state $h$ in two interventions, a gain $\gamma\geq0$, a
coefficient $\beta\geq0$, and $\eta>0$. For an orthogonal projector $P$,
anchor $\mu$, direction $d$ and gate $g\in[0,1]$, the submitted operator is
\begin{align*}
 \alpha&=\min\!\left(1,\frac{\beta\norm{(I-P)(h-\mu)}_2}{\norm h_2+\eta}\right),\\
 F(h)&=h+\gamma(1-g+g\alpha)d+g\alpha P(\mu-h).
\end{align*}
Let starred quantities define a reference operator at the same $h$, and write
$\delta_d=\norm{d-d^*}_2$, $\delta_\mu=\norm{\mu-\mu^*}_2$,
$\delta_P=\norm{P-P^*}_2$, $\delta_g=|g-g^*|$,
$R=\norm{\mu^*-h}_2$, $B=\delta_\mu+R\delta_P$, and
$L_\alpha=\beta/(\norm h_2+\eta)$.
\begin{proposition}[Fixed-state hybrid perturbation]
\begin{equation}\label{eq:hybrid}
 \norm{F(h)-F^*(h)}_2\leq
 \gamma\delta_d+B+(\gamma\norm{d^*}_2+R)
                    (\delta_g+L_\alpha B).
\end{equation}
\end{proposition}
\begin{proof}
Set $v=P(\mu-h)$ and $v^*=P^*(\mu^*-h)$. Projector norms are at most one, so
$\norm{v-v^*}_2\leq B$ and $\norm{v^*}_2\leq R$. The same decomposition
bounds the difference of $(I-P)(h-\mu)$ and its reference by $B$.
Norm and scalar clipping are 1-Lipschitz; therefore
$|\alpha-\alpha^*|\leq L_\alpha B$. Both $g\alpha$ and $1-g+g\alpha$
lie in $[0,1]$, and each changes by at most
$\delta_g+|\alpha-\alpha^*|$. Bounding the direction and correction terms
separately yields \eqref{eq:hybrid}.
\end{proof}

The implementation uses a shared gate in stacked layer space, whose stacked
anchor $\bar\mu$ differs from the individual layer's anchor. For
\[
 g=\max\!\left(0,\frac{z^\top\bar\mu}
                  {\norm z_2\norm{\bar\mu}_2+\epsilon_g}\right),
 \qquad\epsilon_g\geq0,
\]
if both compared $z$ norms are at least $a>0$ and both anchor norms at least
$b>0$, then
\[
 \delta_g\leq\min\!\left(1,
      2\norm{z-z^*}_2/a+2\norm{\bar\mu-\bar\mu^*}_2/b\right).
\]
Change $z$ with $\bar\mu$ fixed, then $\bar\mu$ with $z$ fixed. In each
step the normalized map is $v/(\norm v_2+c)$ for a fixed $c\geq0$, and its
two-endpoint difference is at most $2\norm{v-v^*}_2/\min(\norm v_2,\norm{v^*}_2)$.
The positive-part map is 1-Lipschitz. This proves the stated gate bound,
including numerical denominator regularization.

For a fixed-context score whose gradient norm is at most $K$ along the segment
joining $F(h)$ and $F^*(h)$, the score difference is at most $K$ times
\eqref{eq:hybrid}. Norm-threshold routing at $\tau$ is unchanged whenever
$\norm{z-z^*}_2<|\norm{z^*}_2-\tau|$. These are explicit smooth-score and
routing-margin conditions, not an unconditional generation bound.

\section{Geometry of the additional contextual operator}

The development-selected contextual variant has the different update
\[
 F(h)=h+\gamma d+aP(\mu-h-\gamma d),\qquad a=cg\in[0,1].
\]
Here $P$ is the selected-span projector, $\mu$ the weighted positive-context
anchor, and $d$ the direction from the demonstrations. The uncertainty coefficient
$c$ and cosine gate $g$ are fixed after task adaptation. Set $b=h+\gamma d$. Then
\[
 (I-P)F(h)=(I-P)b,\qquad PF(h)=(1-a)Pb+aP\mu.
\]
Thus the component orthogonal to the selected span is preserved exactly, and
the selected component interpolates toward a composed contextual anchor.
For a fixed task operator,
\begin{align*}
 \norm{F(h)-F(h')}_2^2
 &=\norm{(I-P)(h-h')}_2^2+(1-a)^2\norm{P(h-h')}_2^2\\
 &\leq\norm{h-h'}_2^2.
\end{align*}
This follows from $P^2=P$ and orthogonality of projected and complementary
components. For two fixed operators, put
$R=\norm{\mu^*-h-\gamma d^*}_2$. Changing direction first, then anchor,
then the coefficient/projector gives
\[
 \norm{F(h)-F^*(h)}_2\leq
 \gamma\delta_d+\delta_\mu+R(\delta_P+\delta_a).
\]
The first change uses $\norm{I-aP}_2\leq1$, the second uses
$\norm{aP}_2\leq1$, and the third uses
$\norm{aP-a^*P^*}_2\leq\delta_P+\delta_a$.

The further weighted-skill correction uses
$M=\sum_t w_tP_t$, $b_\mu=\sum_t w_tP_t\mu_t$, with $w_t\geq0$ and
$\sum_tw_t=1$, and applies
\[
 F(h)=h+\gamma d+a\{b_\mu-M(h+\gamma d)\}.
\]
Since $0\preceq M\preceq I$ and $a\in[0,1]$,
$\norm{I-aM}_2\leq1$: this fixed-state map is also nonexpansive.
Rank-truncated weighted bases instead produce an orthogonal projector and
obey the preceding projector argument. These facts explain what is composed;
they do not assert sequential execution of named skills, nonexpansiveness of
the subsequent transformer, or an increase in greedy exact-match accuracy.

\paragraph{Reproduction.}
The machine-readable evidence includes the original coherence audit, all
support-specific bounds, 42 planted recovery checks, 1,200 submitted-hybrid
perturbation checks, 1,000 contextual-geometry checks, and 480 weighted/rank
nonexpansion checks. The last group also includes 120 explicit weighted-sum
identity checks. These are numerical implementation checks rather than substitutes
for the proofs above. All scripts, settings and raw results remain in the
separate rebuttal and optimization directories.
\end{document}
